\documentclass[11pt,a4paper]{article}
\usepackage[T1]{fontenc}
\usepackage{lmodern,amsmath,amssymb,amsthm,mathtools,booktabs,microtype}
\usepackage[margin=27mm]{geometry}
\usepackage[colorlinks=true,linkcolor=blue,citecolor=blue,urlcolor=blue]{hyperref}
\newtheorem{theorem}{Theorem}[section]
\newtheorem{proposition}[theorem]{Proposition}
\newtheorem{lemma}[theorem]{Lemma}
\newtheorem{corollary}[theorem]{Corollary}
\newcommand{\Q}{\mathbb Q}\newcommand{\Z}{\mathbb Z}
\DeclareMathOperator{\Odd}{Odd}
\setlength{\emergencystretch}{2em}
\title{Odd Radicals and Products of Six Linear Factors}
\author{}\date{10 November 2026}
\begin{document}\maketitle
\begin{abstract}
For a fixed product of six rational linear factors, we exhibit three
obstructions to bounding a multiplicative scale by the odd radical of
the product at two arguments. A rational involution gives a negative
square for every prescribed rational scale greater than four. A second
involution gives positive examples of unbounded scale whose first
argument is a sum of two rational squares. A partial symmetry yields
an elliptic family with both arguments positive and the first greater
than one. We determine all rational M\"obius transformations preserving
the six factors and give exact certificates for the identities and
the elliptic non-torsion calculation. The results concern height bounds
for the product; they do not assert solutions of the auxiliary elliptic
and pole equations in a Diophantine definition.
\end{abstract}

\section{Definitions and the height question}

For a nonzero rational number $q$, let $v_p(q)$ be its normalized
$p$-adic valuation. For a finite set $S$ of rational primes, put
\[
 \Odd_S(q)=\prod_{\substack{p\notin S\\v_p(q)\text{ odd}}}p.
\]
Negative valuations are included in this definition. In particular
$\Odd_S(cq^2)=\Odd_S(c)$ for $c,q\in\Q^\times$.
If $x=u/v$ in lowest terms with $v>0$, write
$h(x)=\log\max(|u|,v)$.

Set
\begin{align}
 \Theta&=\left\{0,-\frac12,-\frac14,-\frac25,-1,-\frac13\right\},
       \label{eq:contacts}\\
 F(t)&=\prod_{\theta\in\Theta}(t-\theta)
      =\frac{t(t+1)(2t+1)(3t+1)(4t+1)(5t+2)}{120},
       \label{eq:polynomial}\\
 R(t)&=\frac{F(t)}{(t-1)^8}.\label{eq:R}
\end{align}
An admissible pair $(a,\lambda)$ consists of nonzero rational numbers
such that both $\lambda$ and $a\lambda$ lie outside
$\Theta\cup\{1\}$. We examine bounds of the form
\begin{equation}\label{eq:height-question}
 h(a)\le C\log\Odd_S(R(\lambda)R(a\lambda))+C'
\end{equation}
with constants independent of the admissible pair. The denominator
in~\eqref{eq:R} is a square, but we retain it to record all poles
and to match the rational functions used in the six-contact
construction of~\cite{Conditional}.

Theorems~\ref{thm:universal}, \ref{thm:positive} and
\ref{thm:elliptic} give families for which the radical in
\eqref{eq:height-question} is constant and $h(a)$ is unbounded.
Their domains differ: the first permits every prescribed scale $a>4$,
the second imposes a Gaussian norm condition on $\lambda$, and the
third imposes $\lambda>1$. The additional conditions are stated
separately throughout.

\section{Cancellation for every prescribed scale}

Define
\[
                         \psi(t)=-\frac{t}{5t+1}.
\]
This involution fixes $0,-2/5$ and exchanges $-1/2$ with $-1/3$
and $-1/4$ with $-1$.

\begin{lemma}\label{lem:negative-identity}
As identities of rational functions,
\begin{align}
 F(\psi(t))&=-\frac{F(t)}{(5t+1)^6},\label{eq:psiF}\\
 R(t)R(\psi(t))
 &=-\left(\frac{F(t)(5t+1)}{(t-1)^4(6t+1)^4}\right)^2.
       \label{eq:psiR}
\end{align}
\end{lemma}
\begin{proof}
Put $N=-t$ and $D=5t+1$. The five remaining numerator factors
in $120F(N/D)$ are
\[
 N+D=4t+1,\quad 2N+D=3t+1,\quad 3N+D=2t+1,
\]
\[
                    4N+D=t+1,\qquad 5N+2D=5t+2.
\]
Their product with $N$ is $-120F(t)$, which proves
\eqref{eq:psiF}. Since $\psi(t)-1=-(6t+1)/(5t+1)$,
substitution into~\eqref{eq:R} gives~\eqref{eq:psiR}.
\end{proof}

\begin{theorem}\label{thm:universal}
For every rational $a>4$, set
\[
                        \lambda=-\frac{a+1}{5a}.
\]
Then $(a,\lambda)$ is admissible, $a\lambda=\psi(\lambda)$,
and
\begin{equation}\label{eq:universal-square}
 R(\lambda)R(a\lambda)=
 -\left(
 \frac{25a(a-4)(a-1)(a+1)(2a-3)(3a-2)(4a-1)}
      {24(a+6)^4(6a+1)^4}\right)^2.
\end{equation}
Every finite valuation of this nonzero product is even. Thus
$\Odd_S(R(\lambda)R(a\lambda))=1$ for every finite set $S$.
\end{theorem}
\begin{proof}
One has
\[
 5\lambda+1=-1/a,\qquad
 \psi(\lambda)=-(a+1)/5=a\lambda.
\]
The inequalities
\[
                 -1/4<\lambda<-1/5,\qquad a\lambda<-1
\]
exclude all six contacts and the pole $1$. Moreover
\[
 \lambda-1=-\frac{6a+1}{5a},\qquad
 6\lambda+1=-\frac{a+6}{5a},
\]
so the transformed denominators do not vanish.
Lemma~\ref{lem:negative-identity} proves the negative-square assertion.
Substitution and cancellation yield~\eqref{eq:universal-square},
whose displayed factors are nonzero for $a>4$.
\end{proof}

\begin{corollary}\label{cor:eighth}
Fix rational numbers $B\ge3$ and $T>0$. The inequality
\eqref{eq:height-question} cannot hold for every admissible pair
with
\[
                   a=(2X^2+B)^8+T^8,\qquad X\in\Q,
\]
for any constants $C,C'$ and any fixed finite set $S$.
\end{corollary}
\begin{proof}
Each displayed scale exceeds four, so
Theorem~\ref{thm:universal} gives an admissible slope with product
odd radical one. Taking positive integers $X$ tending to infinity
makes $a$ and hence $h(a)$ tend to infinity.
\end{proof}

In particular the restriction $a=\beta^8+T(n)^8$ in the
weak-$abc$ construction does not, by itself, prevent cancellation:
for a fixed nonzero elliptic index $n$, its trace $T(n)$ is a
fixed positive rational number and $\beta=2X^2+B$. The product
in~\eqref{eq:universal-square} is negative. The theorem asserts
neither the auxiliary local tests nor solvability of a pole predicate.

\section{Contact automorphisms and positive norm slopes}

\begin{proposition}\label{prop:automorphisms}
There are exactly twelve rational M\"obius transformations preserving
$\Theta$. They are listed below. A row $(a,b,c,d;k)$ denotes
\[
 \gamma(t)=\frac{at+b}{ct+d},\qquad
                 F(\gamma(t))(ct+d)^6=kF(t).
\]
\begin{center}
\begin{tabular}{rrrrr}
\toprule
$a$&$b$&$c$&$d$&$k$\\
\midrule
1&0&-5&-1&-1\\
1&0&0&1&1\\
1&1&-7&-4&-27\\
1&1&2&-1&27\\
2&1&-7&-3&1\\
2&1&-3&-2&-1\\
3&1&-8&-3&-1\\
3&1&-7&-2&1\\
4&1&-13&-4&27\\
4&1&-7&-1&-27\\
5&2&-14&-5&-27\\
5&2&-11&-5&27\\
\bottomrule
\end{tabular}
\end{center}
Apart from the identity, the only transformation sending some positive
rational number to a positive number is
\[
                           \tau(t)=\frac{t+1}{2t-1}.
\]
Its positive domain is $t>1/2$. For $t>1$ it satisfies
$0<\tau(t)/t<2$.
\end{proposition}
\begin{proof}
A M\"obius transformation is determined by its values at the three
distinct contacts $0,-1/2,-1/4$. There are exactly
$6\cdot5\cdot4=120$ possible ordered triples of distinct images
in $\Theta$. For a source triple $(x_i)$ and target triple $(y_i)$,
its coefficient vector spans the kernel of
\[
 \begin{pmatrix}
 x_1&1&-x_1y_1&-y_1\\
 x_2&1&-x_2y_2&-y_2\\
 x_3&1&-x_3y_3&-y_3
 \end{pmatrix}.
\]
The signed $3\times3$ minors give that vector explicitly. Testing
the remaining three contacts leaves exactly the twelve rows of the
table. This is a finite rational calculation; the accompanying checker
performs all 120 cases and verifies each polynomial identity.

In each row other than the identity and $\tau$, the numerator
coefficients are positive or zero and the denominator coefficients
are negative. Such a transformation sends every $t>0$ to a negative
number. The asserted positive domain of $\tau$ follows directly
from its denominator. Finally, for $t>1$, the inequality
$\tau(t)/t<2$ is equivalent to
\[
                     4t^2-3t-1=(t-1)(4t+1)>0.
\]
\end{proof}

The rational function identity for $\tau$ in the table gives
\begin{equation}\label{eq:tauR}
 R(t)R(\tau(t))
 =3\left(\frac{3F(t)(2t-1)}{(t-1)^4(t-2)^4}\right)^2.
\end{equation}
Here $\tau(t)-1=(2-t)/(2t-1)$ explains the additional denominator.

\begin{theorem}\label{thm:positive}
For every integer $m\ge2$, set
\[
 \lambda_m=\frac{m^2+1}{2m^2},\qquad
 \mu_m=\frac{3m^2+1}{2},\qquad
 a_m=\frac{m^2(3m^2+1)}{m^2+1}.
\]
Then $(a_m,\lambda_m)$ is admissible, $a_m>4$,
$\mu_m=a_m\lambda_m=\tau(\lambda_m)$, and $a_m\to\infty$.
Both arguments are positive and
\begin{equation}\label{eq:norm}
 \lambda_m=\left(\frac{m+1}{2m}\right)^2+
              \left(\frac{m-1}{2m}\right)^2.
\end{equation}
Moreover $R(\lambda_m)R(a_m\lambda_m)\in3\Q^{\times2}$.
Thus~\eqref{eq:height-question} fails even with positive Gaussian
norm slopes when $3\in S$.
\end{theorem}
\begin{proof}
All relations are immediate substitutions. In particular
\[
 1/2<\lambda_m\le5/8,\qquad \mu_m\ge13/2,
\]
so no argument is a contact or pole. Also
$a_m=3m^2-2+2/(m^2+1)>4$ and $a_m\to\infty$.
Equation~\eqref{eq:norm} proves the norm condition, and
\eqref{eq:tauR} proves the product squareclass assertion.
\end{proof}

For a prescribed rational scale $a>0$, the equation
$\tau(\lambda)=a\lambda$ is
\[
                   2a\lambda^2-(a+1)\lambda-1=0.
\]
It has a rational root precisely when $a^2+10a+1$ is a rational
square. Theorem~\ref{thm:positive} does not impose the further
condition $a=\beta^8+T(n)^8$. Also, Proposition~\ref{prop:automorphisms}
shows that no nonidentity contact automorphism gives
$\gamma(\lambda)=a\lambda$ with $\lambda>1$ and $a>4$.
That observation concerns contact automorphisms only; the next section
gives an obstruction of a different kind.

\section{An elliptic family above one}

Put
\[
 \gamma(t)=\frac1{2(t-1)},\qquad
 H(t)=(2t-1)(3t+1)(5t+2).
\]
Direct multiplication gives
\begin{equation}\label{eq:elliptic-product}
 F(t)F(\gamma(t))=
 \left(\frac{t(t+1)(2t+1)(4t+1)}{480(t-1)^3}\right)^2H(t).
\end{equation}
Thus a rational point on $7y^2=H(t)$, away from the excluded
arguments, gives $R(t)R(\gamma(t))\in7\Q^{\times2}$.

\begin{lemma}\label{lem:elliptic-point}
The elliptic curve
\[
                \mathcal E:\quad Y^2=(X-105)(X+70)(X+84)
\]
contains the non-torsion point $P=(210,2940)$. This point belongs
to the identity component of $\mathcal E(\mathbb R)$.
\end{lemma}
\begin{proof}
Substitution verifies that $P$ lies on the curve. The three roots
are distinct modulo $11$ and modulo $13$, so both primes are of
good reduction. Exact enumeration and the group law give
\[
 \#\mathcal E(\mathbb F_{11})=8,\qquad
                    \operatorname{ord}(P\bmod13)=10.
\]
For clarity the successive nonidentity multiples modulo $13$ are
\[
 (2,2),(11,4),(0,10),(4,6),(1,0),
 (4,7),(0,3),(11,9),(2,11).
\]
The next multiple is the identity. If $P$ were torsion, its order
would therefore be divisible by five. Rational torsion of order
prime to a good reduction prime injects under reduction
\cite[Chapter~VII, Section~3]{Silverman}. At $11$ this would give
an element of order five in a group of order eight, a contradiction.

The largest real root of the cubic is $105$, and $X(P)=210>105$.
Hence $P$ is on the unbounded real branch, which contains the point
at infinity and is the identity component.
\end{proof}

\begin{theorem}\label{thm:elliptic}
There are rational sequences $\lambda_j>1$ and $a_j>4$ with
$a_j\to\infty$ such that $(a_j,\lambda_j)$ is admissible and
\[
                      R(\lambda_j)R(a_j\lambda_j)
                                      \in7\Q^{\times2}.
\]
Consequently~\eqref{eq:height-question} fails on the domain
$\lambda>1$, $a>4$, whenever $7\in S$.
\end{theorem}
\begin{proof}
The change of variables $X=210t$, $Y=1470y$ identifies
$7y^2=H(t)$ with $\mathcal E$. Indeed
\[
 (210t-105)(210t+70)(210t+84)=308700H(t),
 \qquad1470^2=7\cdot308700.
\]
The point $P$ corresponds to $(t,y)=(1,2)$.

The identity component of the real points of an elliptic curve is a
circle group. The multiples of a non-torsion point are dense in that
component: after identification with $\mathbb R/\mathbb Z$, such a
point has irrational coordinate, and its cyclic subgroup is dense.
By Lemma~\ref{lem:elliptic-point}, choose rational multiples $P_j$
tending to $P$ from the local side where $X>210$. This side exists
because $Y(P)\ne0$, so $X$ is a local real coordinate at $P$.
Then
\[
                    \lambda_j=X(P_j)/210>1,\qquad
                    \lambda_j\longrightarrow1.
\]
Set
\[
 \mu_j=\frac1{2(\lambda_j-1)},\qquad
 a_j=\frac{\mu_j}{\lambda_j}
            =\frac1{2\lambda_j(\lambda_j-1)}.
\]
For all sufficiently large $j$, $\mu_j>4\lambda_j>4$ and the
arguments avoid every contact and pole. Equation~\eqref{eq:elliptic-product}
proves the squareclass assertion. Finally $a_j\to\infty$, and
therefore $h(a_j)\to\infty$.
\end{proof}

No assertion in Theorem~\ref{thm:elliptic} imposes that
$\lambda_j-1$ be a sum of two rational squares or that
$a_j=\beta^8+T(n)^8$. Adding either condition requires a further
argument. In particular, this theorem does not settle the product
height question on their intersection.

\section{Quadratic twists and exact verification}

The geometric meaning of the cancellation is straightforward.
For a squarefree polynomial $f$, the equation
\[
                            z^2=f(x)f(y)
\]
records pairs of points on the same quadratic twist of
$w^2=f(x)$. Allowing a fixed nonzero factor on the right records
pairs of twists with a fixed ratio of squareclasses. This pair
surface, including its rational curves and its relation to uniform
height questions, is studied in~\cite[Sections~7--8]{Granville}.
Finiteness of rational points on each fixed genus-two twist does
not give a bound uniform in a varying squareclass. Consequently,
excluding contact automorphism families alone does not establish
\eqref{eq:height-question}.

All identities and finite computations in this paper are verified by
\nolinkurl{check_product_height_obstructions_v19.py} in
\nolinkurl{computation/pole_parameters}. The program uses only rational
arithmetic and the Python standard library. It checks the two
homogeneous sextic identities, the explicit square root in
\eqref{eq:universal-square}, the complete 120-case automorphism
enumeration, and the cleared polynomial identity
\eqref{eq:elliptic-product}. It also verifies the cubic change of
variables, enumerates the groups over $\mathbb F_{11}$ and
$\mathbb F_{13}$, and checks the stated point orders. Its certificate
records these finite results. The argument giving the infinite
elliptic sequence is the real-density proof of
Theorem~\ref{thm:elliptic}.

\begin{thebibliography}{9}\small
\bibitem{Conditional}
J.~Brox and contributors, \emph{Thirteen Variables from a Weak $abc$
Inequality},
\texttt{conditional13\_20261109\_v16.tex},
\url{https://github.com/JoseBrox/diophantine-z-q_improvements}.
\bibitem{Granville}
A.~Granville, \emph{Rational and integral points on quadratic twists
of a given hyperelliptic curve}, International Mathematics Research
Notices \textbf{2007}, article rnm027,
\url{https://doi.org/10.1093/imrn/rnm027}.
\bibitem{Silverman}
J.~H.~Silverman, \emph{The Arithmetic of Elliptic Curves}, second
edition, Graduate Texts in Mathematics \textbf{106}, Springer, 2009,
\url{https://doi.org/10.1007/978-0-387-09494-6}.
\end{thebibliography}
\end{document}
